Em Cunha, p. 338: \enquote{a agitação começou em abril de 2003, quando uma carta assinada por índios katukina do Acre chegou ao Ministério do Meio Ambiente. A carta afirmava que o uso da secreção de certa rã arborícola (como a zoologia chama uma perereca), difundido em várias cidades do país havia alguns anos, derivava do conhecimento tradicional katukina, e que este estava sendo indevidamente apropriado. A então ministra Marina Silva, que como se sabe é acriana e filha de seringueiros, se comprometeu a fazer desse caso um exemplo positivo de defesa de direitos sobre conhecimentos tradicionais (...). Fui então chamada pelo Ministério a participar de um grupo de trabalho para examinar o caso, mas logo fiz notar que ele envolvia questões complexas e portanto não era particularmente promissor para estabelecer o desejado paradigma positivo.}